\documentclass[12pt]{article}
\usepackage[utf8]{inputenc}
\usepackage[spanish]{babel}
\usepackage{listings}
\usepackage{xcolor}
\usepackage{amsmath}
\usepackage{amssymb}
\usepackage{geometry}
\geometry{margin=1in}

% Configuración de colores para código
\definecolor{codegreen}{rgb}{0,0.6,0}
\definecolor{codegray}{rgb}{0.5,0.5,0.5}
\definecolor{codepurple}{rgb}{0.58,0,0.82}
\definecolor{backcolour}{rgb}{0.95,0.95,0.92}

\lstdefinestyle{haskellstyle}{
    backgroundcolor=\color{backcolour},
    commentstyle=\color{codegreen},
    keywordstyle=\color{magenta},
    numberstyle=\tiny\color{codegray},
    stringstyle=\color{codepurple},
    basicstyle=\ttfamily\footnotesize,
    breakatwhitespace=false,
    breaklines=true,
    captionpos=b,
    keepspaces=true,
    numbers=left,
    numbersep=5pt,
    showspaces=false,
    showstringspaces=false,
    showtabs=false,
    tabsize=2
}

\lstset{style=haskellstyle}

\title{Combinador Y: Recursión sin Auto-Referencia\\
\large Diálogo para Presentación}
\author{}
\date{}

\begin{document}

\maketitle

\section*{Guión de Presentación}

\subsection*{Introducción (Plantear la pregunta)}

\begin{quote}
\textbf{Presentador:} ``Hoy vamos a responder una pregunta que puede parecer obvia al principio: ¿Cuál es la diferencia entre un intérprete normal y uno que usa el combinador Y?''

\textbf{Presentador:} ``A simple vista, ambos interpretan expresiones aritméticas de la misma forma. Veamos un ejemplo...''
\end{quote}

\subsection*{Demostración 1: Mismo comportamiento}

\begin{quote}
\textbf{Presentador:} ``Evaluemos la expresión \texttt{(+ (* 3 4) 5)} con ambos enfoques:''
\end{quote}

\begin{verbatim}
Evaluando: (+ (* 3 4) 5)
Con recursión tradicional: 17
Con combinador Y:          17

Resultado: ¡Son iguales!
\end{verbatim}

\begin{quote}
\textbf{Presentador:} ``Como vemos, el \textit{comportamiento} es idéntico. Pero entonces, ¿dónde está la diferencia?''

\textbf{Presentador:} ``La diferencia está en \textit{cómo implementamos la recursión}. Veamos el código...''
\end{quote}

\subsection*{Demostración 2: Código lado a lado}

\begin{quote}
\textbf{Presentador:} ``Primero, el enfoque tradicional:''
\end{quote}

\begin{lstlisting}[language=Haskell, caption=Recursión Tradicional]
evalTradicional :: ASAValues -> Int
evalTradicional expr = case expr of
    NumV n -> n

    -- La función SE LLAMA A SÍ MISMA por nombre
    AddV e1 e2 -> evalTradicional e1 + evalTradicional e2
    --            ^^^^^^^^^^^^^^^ auto-referencia explicita

    SubV e1 e2 -> evalTradicional e1 - evalTradicional e2
    MultV e1 e2 -> evalTradicional e1 * evalTradicional e2
    DivV e1 e2 -> evalTradicional e1 `div` evalTradicional e2
\end{lstlisting}

\begin{quote}
\textbf{Presentador:} ``Noten cómo \texttt{evalTradicional} se llama a sí misma directamente por nombre. Esto es recursión explícita.''

\textbf{Presentador:} ``Ahora veamos el enfoque con combinador Y:''
\end{quote}

\begin{lstlisting}[language=Haskell, caption=Con Combinador Y]
evalMaker :: (ASAValues -> Int) -> (ASAValues -> Int)
evalMaker evalRec expr = case expr of
    NumV n -> n

    -- NO se llama a si misma, usa el parametro 'evalRec'
    AddV e1 e2 -> evalRec e1 + evalRec e2
    --            ^^^^^^^ NO es auto-referencia, es un parametro

    SubV e1 e2 -> evalRec e1 - evalRec e2
    MultV e1 e2 -> evalRec e1 * evalRec e2
    DivV e1 e2 -> evalRec e1 `div` evalRec e2

-- El combinador Y "inyecta" la recursion
evalExpr :: ASAValues -> Int
evalExpr = yFix evalMaker
\end{lstlisting}

\begin{quote}
\textbf{Presentador:} ``¡Aquí está la diferencia clave! \texttt{evalMaker} \textbf{NUNCA} se llama a sí misma por nombre. En lugar de eso, recibe la recursión como un parámetro \texttt{evalRec}.''

\textbf{Presentador:} ``El combinador Y es quien 'inyecta' la recursión desde afuera.''
\end{quote}

\newpage
\subsection*{Explicación Técnica: ¿Cómo funciona?}

\begin{quote}
\textbf{Presentador:} ``Veamos paso a paso cómo el combinador Y hace esta 'magia'.''
\end{quote}

\subsubsection*{Definición del Combinador Y}

\begin{lstlisting}[language=Haskell]
yFix :: ((a -> b) -> (a -> b)) -> (a -> b)
yFix f = f (yFix f)
\end{lstlisting}

\begin{quote}
\textbf{Presentador:} ``Esta definición tan simple tiene una propiedad poderosa: encuentra el punto fijo de una función.''
\end{quote}

\subsubsection*{Propiedad de Punto Fijo}

\begin{quote}
\textbf{Presentador:} ``Un punto fijo de una función $f$ es un valor $x$ tal que $f(x) = x$.''

\textbf{Presentador:} ``Para nuestro caso:''
\end{quote}

\begin{equation*}
\texttt{evalExpr} = \texttt{yFix evalMaker} = \texttt{evalMaker evalExpr}
\end{equation*}

\begin{quote}
\textbf{Presentador:} ``Es decir, \texttt{evalExpr} es el punto fijo de \texttt{evalMaker}.''
\end{quote}

\subsubsection*{Expansión Paso a Paso}

\begin{quote}
\textbf{Presentador:} ``Veamos cómo se expande la evaluación de \texttt{(+ 2 3)}:''
\end{quote}

\begin{align*}
&\texttt{evalExpr (AddV (NumV 2) (NumV 3))} \\
&= \texttt{(yFix evalMaker) (AddV (NumV 2) (NumV 3))} \\
&= \texttt{(evalMaker (yFix evalMaker)) (AddV (NumV 2) (NumV 3))} \quad \text{(def. de yFix)} \\
&= \texttt{(evalMaker evalExpr) (AddV (NumV 2) (NumV 3))} \quad \text{(punto fijo)} \\
&= \texttt{evalExpr (NumV 2) + evalExpr (NumV 3)} \quad \text{(expansión de evalMaker)} \\
&= \texttt{2 + 3} \\
&= \texttt{5}
\end{align*}

\begin{quote}
\textbf{Presentador:} ``¡La recursión funciona sin que \texttt{evalMaker} se llame a sí misma!''
\end{quote}

\newpage
\subsection*{Ejemplo con Factorial}

\begin{quote}
\textbf{Presentador:} ``Para entenderlo mejor, veamos un ejemplo más familiar: el factorial.''
\end{quote}

\begin{verbatim}
=== COMPARACIÓN: Factorial Tradicional vs Con Y ===
factorialTradicional 5 = 120
factorial (con Y)      5 = 120
\end{verbatim}

\begin{lstlisting}[language=Haskell, caption=Factorial Tradicional]
factorialTradicional :: Int -> Int
factorialTradicional n
    | n <= 0    = 1
    | otherwise = n * factorialTradicional (n - 1)
                      -- ^^^^^^^^^^^^^^^^^^^^ auto-referencia
\end{lstlisting}

\begin{lstlisting}[language=Haskell, caption=Factorial con Y]
factorialMaker :: (Int -> Int) -> (Int -> Int)
factorialMaker rec n
    | n <= 0    = 1
    | otherwise = n * rec (n - 1)
                      -- ^^^ parametro, no auto-referencia

factorial :: Int -> Int
factorial = yFix factorialMaker
\end{lstlisting}

\subsection*{Diferencias Resumidas}

\begin{quote}
\textbf{Presentador:} ``Resumamos las diferencias clave:''
\end{quote}

\begin{enumerate}
\item \textbf{Recursión Tradicional:}
    \begin{itemize}
    \item La función se llama a sí misma por nombre
    \item Requiere que el lenguaje soporte recursión explícita
    \item Es lo que normalmente escribimos
    \end{itemize}

\item \textbf{Combinador Y:}
    \begin{itemize}
    \item La función \textit{nunca} se refiere a sí misma
    \item La recursión viene ``desde afuera'' vía el combinador Y
    \item Demuestra que la recursión puede implementarse sin auto-referencia
    \end{itemize}
\end{enumerate}

\newpage
\subsection*{Importancia Teórica}

\begin{quote}
\textbf{Presentador:} ``Pero, ¿por qué es importante esto?''
\end{quote}

\subsubsection*{1. Base del Cálculo Lambda}

\begin{quote}
\textbf{Presentador:} ``El Cálculo Lambda puro \textbf{no tiene recursión explícita}. No puedes nombrar funciones para llamarlas recursivamente.''

\textbf{Presentador:} ``El combinador Y demuestra que \textit{no necesitas} recursión como primitiva del lenguaje. Puedes implementarla usando solo:''
\end{quote}

\begin{itemize}
\item Funciones
\item Aplicación de funciones
\item Funciones de orden superior
\end{itemize}

\subsubsection*{2. Recursión sin Auto-Referencia}

\begin{quote}
\textbf{Presentador:} ``Imagina que te prohíben usar el nombre de una función dentro de su propio cuerpo. ¿Cómo harías recursión?''

\textbf{Presentador:} ``El combinador Y es la respuesta: inyecta la recursión desde afuera.''
\end{quote}

\subsubsection*{3. Aplicaciones Prácticas}

\begin{itemize}
\item \textbf{Lenguajes funcionales:} Base teórica para implementar recursión
\item \textbf{Teoría de tipos:} Entender recursión en sistemas de tipos
\item \textbf{Optimización:} Algunos compiladores usan puntos fijos para optimizar
\item \textbf{Semántica:} Definir semántica de lenguajes con recursión
\end{itemize}

\subsection*{Conclusión}

\begin{quote}
\textbf{Presentador:} ``Entonces, ¿cuál es la diferencia entre un intérprete normal y uno con combinador Y?''

\textbf{Presentador:} ``El \textit{comportamiento} es idéntico, pero la \textit{implementación} es radicalmente diferente:''
\end{quote}

\begin{center}
\begin{tabular}{|l|l|}
\hline
\textbf{Aspecto} & \textbf{Diferencia} \\
\hline
Comportamiento & Idéntico \\
Auto-referencia & Tradicional: Sí, Combinador Y: No \\
Recursión & Tradicional: Explícita, Combinador Y: Inyectada \\
Requisito del lenguaje & Tradicional: Recursión nativa, Combinador Y: Solo HOF \\
Importancia & Tradicional: Práctica, Combinador Y: Teórica y práctica \\
\hline
\end{tabular}
\end{center}

\begin{quote}
\textbf{Presentador:} ``El combinador Y demuestra un resultado profundo: \textbf{la recursión no es una primitiva necesaria}, sino que puede ser \textit{derivada} de funciones de orden superior.''

\textbf{Presentador:} ``Esta es la base teórica del Cálculo Lambda y de lenguajes funcionales modernos.''
\end{quote}

\newpage
\section*{Autointerpretación y Metaprogramación con el Combinador Y}

\subsection*{La Conexión Clave}

\begin{quote}
\textbf{Presentador:} ``Ahora llegamos a la parte más interesante: ¿Cómo el combinador Y permite la autointerpretación y la metaprogramación?''

\textbf{Presentador:} ``La respuesta está en que el combinador Y permite que un intérprete \textit{se pase a sí mismo como argumento}.''
\end{quote}

\subsubsection*{¿Qué es la Autointerpretación?}

\begin{quote}
\textbf{Presentador:} ``Un autointérprete es un programa que puede interpretar su propio lenguaje. Es decir:''
\end{quote}

\begin{itemize}
\item El intérprete está escrito en el lenguaje $L$
\item Interpreta programas también escritos en $L$
\item $L_{\text{anfitrión}} = L_{\text{objeto}}$
\end{itemize}

\begin{quote}
\textbf{Presentador:} ``Sin el combinador Y, esto es \textbf{imposible} en cálculo lambda puro, porque no podríamos expresar la recursión necesaria para que el intérprete se llame a sí mismo.''
\end{quote}

\subsubsection*{Implementación: Autointérprete Básico}

\begin{quote}
\textbf{Presentador:} ``Veamos cómo se implementa un autointérprete usando el combinador Y:''
\end{quote}

\begin{lstlisting}[language=Haskell, caption=Autointérprete con Combinador Y]
-- Tipo de datos para representar expresiones del lenguaje
data ASAValues
    = NumV Int
    | AddV ASAValues ASAValues
    | SubV ASAValues ASAValues
    | MultV ASAValues ASAValues
    | DivV ASAValues ASAValues
    deriving (Show, Eq)

-- El "generador" del interprete (funcion casi-recursiva)
-- Recibe el interprete como parametro 'self'
interpreteMaker :: (ASAValues -> Int) -> (ASAValues -> Int)
interpreteMaker self expr = case expr of
    -- Casos base
    NumV n -> n

    -- Casos recursivos: el interprete se usa a si mismo
    -- para evaluar subexpresiones
    AddV e1 e2 -> self e1 + self e2
    --            ^^^^ El interprete se pasa a si mismo

    SubV e1 e2 -> self e1 - self e2
    MultV e1 e2 -> self e1 * self e2
    DivV e1 e2 ->
        let v2 = self e2
        in if v2 == 0
           then error "Division por cero"
           else self e1 `div` v2

-- El AUTOINTERPRETE: se obtiene aplicando Y al generador
-- Esto "cierra el ciclo" de autorreferencia
interprete :: ASAValues -> Int
interprete = yFix interpreteMaker

-- Propiedad clave: interprete = interpreteMaker interprete
-- El interprete se recibe a si mismo como argumento!
\end{lstlisting}

\begin{quote}
\textbf{Presentador:} ``Observen la línea clave: \texttt{interprete = yFix interpreteMaker}''

\textbf{Presentador:} ``Esto garantiza que \texttt{interprete = interpreteMaker interprete}, es decir, el intérprete se pasa a sí mismo como el parámetro \texttt{self}.''

\textbf{Presentador:} ``¡Esto es autointerpretación! El intérprete se usa a sí mismo para evaluar subexpresiones.''
\end{quote}

\newpage
\subsection*{¿Por Qué el Combinador Y Es Necesario?}

\begin{quote}
\textbf{Presentador:} ``Pregunta: ¿Por qué no podemos hacer esto sin el combinador Y?''
\end{quote}

\begin{enumerate}
\item \textbf{En cálculo lambda puro:} No hay nombres, no puedes escribir una función que se llame a sí misma

\item \textbf{Sin combinador Y:} No hay forma de crear el ciclo de autorreferencia necesario

\item \textbf{Con combinador Y:} El punto fijo automáticamente inyecta el intérprete en sí mismo
\end{enumerate}

\begin{quote}
\textbf{Presentador:} ``El combinador Y es el \textit{único} mecanismo que permite autointerpretación en cálculo lambda puro.''
\end{quote}

\subsection*{De Autointerpretación a Metaprogramación}

\begin{quote}
\textbf{Presentador:} ``La autointerpretación abre la puerta a la metaprogramación. ¿Por qué?''
\end{quote}

\subsubsection*{Cadena de Implicaciones}

\begin{align*}
\text{Combinador Y} &\Rightarrow \text{Autointerpretación posible} \\
&\Rightarrow \text{El código puede tratarse como datos} \\
&\Rightarrow \text{El programa puede manipular código} \\
&\Rightarrow \text{Metaprogramación}
\end{align*}

\subsubsection*{Ejemplos de Metaprogramación}

\begin{quote}
\textbf{Presentador:} ``Veamos ejemplos concretos de metaprogramación habilitados por el combinador Y:''
\end{quote}

\textbf{Ejemplo 1: Inspección de Código}

\begin{lstlisting}[language=Haskell, caption=Analizar estructura del programa]
-- El programa puede examinar su propia estructura
contieneDivision :: ASAValues -> Bool
contieneDivision expr = case expr of
    NumV _ -> False
    DivV _ _ -> True  -- Encontro una division!
    AddV e1 e2 -> contieneDivision e1 || contieneDivision e2
    SubV e1 e2 -> contieneDivision e1 || contieneDivision e2
    MultV e1 e2 -> contieneDivision e1 || contieneDivision e2

-- Uso: detectar si una expresion tiene divisiones
-- antes de evaluarla
evaluacionSegura :: ASAValues -> Maybe Int
evaluacionSegura expr =
    if contieneDivision expr
    then Nothing  -- Advertir: contiene division
    else Just (interprete expr)
\end{lstlisting}

\begin{quote}
\textbf{Presentador:} ``Esto es metaprogramación: el código examina código antes de ejecutarlo.''
\end{quote}

\newpage
\textbf{Ejemplo 2: Optimización de Código}

\begin{lstlisting}[language=Haskell, caption=Transformar y optimizar programas]
-- El programa puede reescribirse a si mismo
optimizar :: ASAValues -> ASAValues
optimizar expr = case expr of
    -- Optimizacion: eliminar sumas con cero
    AddV (NumV 0) e -> optimizar e
    AddV e (NumV 0) -> optimizar e

    -- Optimizacion: eliminar multiplicaciones por 1
    MultV (NumV 1) e -> optimizar e
    MultV e (NumV 1) -> optimizar e

    -- Optimizacion: multiplicacion por 0 = 0
    MultV (NumV 0) _ -> NumV 0
    MultV _ (NumV 0) -> NumV 0

    -- Recursivamente optimizar subexpresiones
    AddV e1 e2 -> AddV (optimizar e1) (optimizar e2)
    SubV e1 e2 -> SubV (optimizar e1) (optimizar e2)
    MultV e1 e2 -> MultV (optimizar e1) (optimizar e2)
    DivV e1 e2 -> DivV (optimizar e1) (optimizar e2)

    -- Caso base
    NumV n -> NumV n

-- Evaluar con optimizacion
evalOptimizado :: ASAValues -> Int
evalOptimizado = interprete . optimizar

-- Ejemplo:
-- Original:     (+ (* 0 5) (+ 3 0))
-- Optimizado:   (+ 0 3) -> 3
\end{lstlisting}

\begin{quote}
\textbf{Presentador:} ``El programa se transforma a sí mismo antes de ejecutarse. ¡Esto es metaprogramación en acción!''
\end{quote}

\textbf{Ejemplo 3: Evaluación Parcial}

\begin{lstlisting}[language=Haskell, caption=Evaluacion parcial con conocimiento en tiempo de compilacion]
-- Evaluar parcialmente cuando tengamos valores conocidos
evaluarParcial :: ASAValues -> ASAValues
evaluarParcial expr = case expr of
    -- Si ambos operandos son numeros, evaluar inmediatamente
    AddV (NumV a) (NumV b) -> NumV (a + b)
    SubV (NumV a) (NumV b) -> NumV (a - b)
    MultV (NumV a) (NumV b) -> NumV (a * b)
    DivV (NumV a) (NumV b) | b /= 0 -> NumV (a `div` b)

    -- Si solo uno es numero, optimizar lo posible
    AddV e1 e2 -> AddV (evaluarParcial e1) (evaluarParcial e2)
    SubV e1 e2 -> SubV (evaluarParcial e1) (evaluarParcial e2)
    MultV e1 e2 -> MultV (evaluarParcial e1) (evaluarParcial e2)
    DivV e1 e2 -> DivV (evaluarParcial e1) (evaluarParcial e2)

    NumV n -> NumV n

-- Ejemplo:
-- Entrada:   (+ (+ 2 3) (* 4 5))
-- Parcial:   (+ 5 20)
-- Parcial2:  25
\end{lstlisting}

\subsection*{¿Por Qué Esto Requiere el Combinador Y?}

\begin{quote}
\textbf{Presentador:} ``Estas capacidades de metaprogramación son posibles porque:''
\end{quote}

\begin{enumerate}
\item \textbf{Autointerpretación:} El combinador Y permite que el intérprete se procese a sí mismo

\item \textbf{Código = Datos:} Las expresiones (tipo \texttt{ASAValues}) son datos que podemos manipular

\item \textbf{Recursión sin límites:} El combinador Y permite procesar estructuras arbitrariamente anidadas

\item \textbf{Sin primitivas:} Todo esto funciona en cálculo lambda puro, sin necesidad de características especiales del lenguaje
\end{enumerate}

\subsection*{Demostración: Metaprogramación en Acción}

\begin{lstlisting}[language=Haskell, caption=Pipeline de metaprogramacion completo]
-- Pipeline de transformaciones metaprogramaticas
evaluarConMetaprogramacion :: ASAValues -> Int
evaluarConMetaprogramacion expr =
    let paso1 = optimizar expr           -- Transformacion 1
        paso2 = evaluarParcial paso1     -- Transformacion 2
    in interprete paso2                  -- Interpretacion final

-- Ejemplo completo:
expr = AddV (MultV (NumV 0) (NumV 999))
            (AddV (NumV 2) (NumV 3))

-- Paso 1 (optimizar):     (+ 0 (+ 2 3))
-- Paso 2 (eval parcial):  (+ 0 5)
-- Paso 3 (optimizar):     5
-- Resultado:              5

-- SIN metaprogramacion, habria evaluado: 0*999 + 2 + 3
-- CON metaprogramacion, simplifica antes de evaluar
\end{lstlisting}

\begin{quote}
\textbf{Presentador:} ``El combinador Y no solo permite recursión. Habilita que el lenguaje razone sobre sí mismo, se optimice a sí mismo, y se transforme a sí mismo.''

\textbf{Presentador:} ``Esto es el fundamento de:''
\end{quote}

\begin{itemize}
\item Compiladores optimizadores
\item Macros en Lisp/Scheme
\item Reflection en lenguajes modernos
\item JIT compilation
\item Partial evaluation
\end{itemize}

\begin{quote}
\textbf{Presentador:} ``Todo esto es posible gracias al simple pero poderoso combinador Y.''
\end{quote}

\newpage
\section*{Código Fuente Completo}

\subsection*{YCombinator.hs}

\begin{lstlisting}[language=Haskell]
module YCombinator where

-- Combinador Y (combinador de punto fijo)
-- Version practica para Haskell con evaluacion perezosa
--
-- Y = λf. (λx. f (x x)) (λx. f (x x))
--
-- En Haskell, gracias a la evaluacion perezosa, podemos usar
-- la forma mas simple: yFix f = f (yFix f)
--
-- Esta funcion toma una funcion casi-recursiva y retorna
-- su punto fijo, permitiendo recursion sin autorreferencia.

yFix :: ((a -> b) -> (a -> b)) -> (a -> b)
yFix f = f (yFix f)
\end{lstlisting}

\subsection*{Evaluator.hs (Fragmento relevante)}

\begin{lstlisting}[language=Haskell]
-- Generador de la funcion de evaluacion
-- Esta es una funcion casi recursiva que toma como argumento
-- la funcion recursiva misma (evalRec)
--
-- Nota importante: evalMaker NUNCA se refiere a si misma por nombre
-- La recursion se logra completamente a traves del combinador Y

evalMaker :: (ASAValues -> Int) -> (ASAValues -> Int)
evalMaker evalRec expr = case expr of
    NumV n -> n
    AddV e1 e2 -> evalRec e1 + evalRec e2
    SubV e1 e2 -> evalRec e1 - evalRec e2
    MultV e1 e2 -> evalRec e1 * evalRec e2
    DivV e1 e2 ->
        let v2 = evalRec e2
        in if v2 == 0
           then error "Division por cero"
           else evalRec e1 `div` v2

-- Funcion de evaluacion obtenida aplicando el combinador Y
-- Propiedad de punto fijo: evalExpr = evalMaker evalExpr
evalExpr :: ASAValues -> Int
evalExpr = yFix evalMaker
\end{lstlisting}

\section*{Referencias y Lecturas Adicionales}

\begin{itemize}
\item \textbf{Cálculo Lambda:} Church, Alonzo. ``The Calculi of Lambda-Conversion'' (1941)
\item \textbf{Combinadores:} Curry, Haskell B. ``Combinatory Logic'' (1958)
\item \textbf{Punto Fijo:} Tarski's Fixed Point Theorem
\item \textbf{Implementación práctica:} Pierce, Benjamin C. ``Types and Programming Languages'' (2002)
\end{itemize}

\end{document}
